% Generated draft: requires mathematical and source review.
\documentclass[11pt]{article}
\usepackage[margin=25mm]{geometry}
\usepackage{fontspec}
\setmainfont{Latin Modern Roman}
\usepackage{amsmath,amssymb,mathtools,unicode-math}
\setmathfont{Latin Modern Math}
\usepackage{enumitem,needspace,xurl,hyperref}
\setlength{\parindent}{0pt}
\setlength{\parskip}{5pt}
\setlength{\emergencystretch}{3em}
\newcommand{\N}{\mathbb{N}}
\newcommand{\Z}{\mathbb{Z}}
\newcommand{\Q}{\mathbb{Q}}
\newcommand{\R}{\mathbb{R}}
\newcommand{\C}{\mathbb{C}}
\newcommand{\F}{\mathbb{F}}
\newcommand{\A}{\mathbb{A}}
\newcommand{\PP}{\mathbb{P}}
\newcommand{\E}{\mathbb{E}}
\newcommand{\eps}{\varepsilon}
\newcommand{\dd}{\,\mathrm{d}}
\begin{document}
\section*{UnsolvedMath ID: 1900011}
\textbf{Source problem code:} catalogue entry
\subsection{Definitions and mathematical statement}
The two-dimensional torus is $\mathbb T^2=\mathbb R^2/\mathbb Z^2$. A torus decomposition is conventionally a finite cell decomposition of this surface into vertices, edges, and polygonal two-cells, with their incidence data. The source does not say whether lattice lengths, areas, or other integer-geometric labels are part of the data.

In the geometric sail interpretation of two-dimensional continued fractions, a periodic object carries a group of periods acting on its sail, and a torus decomposition is obtained as a quotient of the sail's faces, edges, and vertices. A rank-two free abelian period group acting freely, properly discontinuously, and cocompactly on a plane-like sail is the usual model for such a quotient. This describes the intended geometric framework conditionally; the record does not define its admissible sails or period groups.

\textbf{Problem (V. Arnold).} Describe exactly the torus decompositions realized by periodic two-dimensional continued fractions. More explicitly, once the source's notions of periodic continued fraction and realization equivalence are fixed, give necessary and sufficient conditions on a torus decomposition $D$ for the existence of such a continued fraction with sail $\mathcal S$ and period group $\Lambda$ whose quotient decomposition $\mathcal S/\Lambda$ is equivalent to $D$. The source asks for a complete classification, not merely examples or a list of necessary conditions. Its supplied research summary reports partial examples but no verified full classification, with moderate-confidence OPEN-TRIAGE status.
\subsection{Short English statement}
Which decompositions of a torus can occur as the periodic quotient of a two-dimensional geometric continued fraction? The exact sail, period, and equivalence conventions are missing from the supplied statement.
\subsection{Sources}
\begin{itemize}
\item O. Karpenkov, \emph{Open problems in geometry of continued fractions} (2017), Problem 11, attributed to V. Arnold, \url{https://arxiv.org/abs/1712.01450}; source-TeX extraction, accessed 2026-07-29.
\item The supplied summary refers to partial results and examples but gives no further publication or theorem citation establishing their scope or resolving the classification.
\item Local registry \texttt{data/unsolvedmath/problems.json}, permanent ID $1900011$, source code catalogue entry, \emph{Geometry of Continued Fractions --- Combinatorial structure of sails}. The record labels this open and OPEN-TRIAGE, with NEEDS\_REVIEW source-status and rights qualifications.
\end{itemize}
\end{document}
